\documentclass[10pt,a4paper]{amsart}
\usepackage[margin=19mm]{geometry}
\usepackage{amsmath,amssymb,mathtools}
\usepackage[scr=rsfs,cal=euler]{mathalfa}
\usepackage{xfrac}
\usepackage[unicode,hidelinks,bookmarks=false]{hyperref}
\newcommand{\ie}{i.e.\ }
\newcommand{\cf}{cf.\ }
\newcommand{\resp}{respectively\ }
\newcommand{\Subsection}{\subsection}
\providecommand{\nobreakdash}{\nobreak\hbox{-}}
\setlength{\emergencystretch}{2em}
\multlinegap0pt
\usepackage{bm}
\usepackage{bbm}
\usepackage{dsfont}
\usepackage{xspace}
\usepackage{cancel}
\usepackage{mathtools}
\usepackage{amsthm}
\makeatletter
\@ifundefined{theorem}{\newtheorem{theorem}{Theorem}}{}
\@ifundefined{corollary}{\newtheorem{corollary}[theorem]{Corollary}}{}
\makeatother
\newcommand{\T}{\mathrm{T}}
\newcommand{\alt}{{\mathrm{alt}}}
\title[Error term in the Cohen-Lenstra heuristic via random matrix ]{Error term in the Cohen-Lenstra heuristic via random matrix approach}
\author{Yue Xu, Xiuwu Zhu}
\date{Source: arXiv:2308.02232v2 (CC BY 4.0)}
\begin{document}
\maketitle

\noindent\emph{Source and licence.} Error term in the Cohen-Lenstra heuristic via random matrix approach. Version arXiv:2308.02232v2 (CC BY 4.0), distributed under the Creative Commons Attribution 4.0 International licence, as recorded in the arXiv licence metadata for this exact version and shown on the abstract page https://arxiv.org/abs/2308.02232v2. Full rights notice: licenses/arxiv-2308.02232-CC-BY-4.0.txt.\par\smallskip

\noindent\textit{Editorial scope.} Counted unit: the Corollary whose source label is \texttt{cor:alt\_converg} in the article (source0.tex lines 1120--1129, source label \texttt{cor:alt\_converg}), delivered together with the article's own complete original proof of it (source0.tex lines 1131--1186, one \texttt{proof} environment of 56 lines). No mathematics is written, repaired or re-derived: statement and proof are the article's own text line for line. Required context retained uncounted: thm:main\_thm\_spec (lines 488--501). Every definition, display and label of those lines is retained unchanged. Omitted from this package and not redistributed: 1--487, 502--1119, 1130--1130, 1187--1388. Nothing inside a retained excerpt is omitted or altered: every retained body is delivered exactly as the article's lines are written, so no omission receipt of any kind accompanies this package. Citations: the retained excerpts carry no citation command at all, so no bibliography entry of the article is transcribed and no reference is redistributed. Reference adaptation: none was needed and none was made; every reference inside the delivered unit resolves inside it, so no unresolved reference or citation is produced. Printed slips kept literal: no printed word, symbol, hypothesis or number of the counted statement or of its proof was altered. Layout and class: the article's own class is \texttt{amsart} with packages including amscd, amsthm, amsfonts, amssymb, amsmath, enumerate, geometry, graphicx, caption, float, subcaption; the wrapper is the lane's pinned \texttt{amsart} editorial wrapper at 10pt on A4, which restates the theorem-like environments the retained text needs and adds the article's own macro definitions that the retained text needs (\texttt{\textbackslash T}, \texttt{\textbackslash alt}), with extra wrapper packages (bm, bbm, dsfont, xspace, cancel, mathtools). The delivered numbering is the unit's own. Assets: no retained span contains a figure environment, a graphics-inclusion command, a PDF-page inclusion, a table or a TikZ picture; no image file is copied into this package, the package declares no asset dependency and no image file is redistributed with it. Licence: version arXiv:2308.02232v2 of this article is distributed under the Creative Commons Attribution 4.0 International licence, as recorded in the arXiv licence metadata for this exact version and shown on the abstract page https://arxiv.org/abs/2308.02232v2. A full rights notice is delivered at licenses/arxiv-2308.02232-CC-BY-4.0.txt. Characters: every retained excerpt is pure ASCII, so the delivered engine, XeLaTeX with its default Latin Modern text font, reports no missing character.
\begin{theorem}\label{thm:main_thm_spec}
Assume further that $P$ is a compact operator on $\ell^2(\pi)$. Let $\lambda_0, \lambda_1, \lambda_2, \cdots$ be all eigenvalues of $P$ with non-increasing absolute value. Then for any $\mu \in \ell^2(\pi)$,
\[
\|\mu P^n - \sum_{i=0}^{k} \lambda_i^n \mu_i\|_{tv} = O(|\lambda_{k+1}|^{n}).
\]
Here, $\mu_i$ is the $\lambda_i$-component in the spectral decomposition of $\mu$, and the implicit constant is less than $\|\mu\|_\pi$. In particular,
\[
\|\mu P^n - \mu_0\|_{tv} = \begin{cases}
    \|\mu_1\|_{tv} \cdot |\lambda_1|^{n} + O(|\lambda_2|^{n}), & \text{if } |\lambda_1| > |\lambda_2|, \\
    (\|\mu_1 + (-1)^n \mu_2\|_{tv}) \cdot |\lambda_1|^{n} + O(|\lambda_3|^{n}), & \text{if } |\lambda_1| = |\lambda_2|.
  \end{cases}
\]
Here, $\mu_0 = (\mu \cdot \mathbf{1}) \pi$, $\mathbf{1} = (1, 1,\cdots, 1, \cdots)^\T$, and the implicit constant does not exceed $\left(\|\mu\|_\pi^2 - (\mu \cdot \mathbf{1})^2\right)^{1/2}$.
\end{theorem}

\begin{corollary}\label{cor:alt_converg}
The convergence rates are:
\[
\|\delta_0 P_\alt^n - \pi_\alt\|_{tv} = \frac{2\alpha(q)}{(q-1)^2(q+1)} q^{-2n} + O(q^{-4n}),
\]
\[
\|\delta_0 Q_\alt^n - \pi_\alt'\|_{tv} = \frac{\alpha(q) q}{(q-1)(q+1)} q^{-2n} + O(q^{-4n}),
\]
with implicit constants less than $(\eta_1(q)^2\alpha(q)^{-2}-1)^{1/2}$ and $(\alpha(q)^{-2}-1)^{1/2}$ respectively.
\end{corollary}

\begin{proof}
Let us first analyze the case for $P_\alt$. We begin by constructing the eigenvector associated with the eigenvalue $q^{-2}$:
\[
\nu := \pi_\alt - \tfrac{q}{q+1}(\pi_\alt \circ q^2) \in V_{q^{-2}}.
\]
Furthermore, we observe that the following combination belongs to the eigenspace $V_{q^{-4}}$:
\[
(1+q^2)(1+q^{-1})(\pi_\alt-\pi_\alt\circ q^2)+\pi_\alt\circ q^4 \in V_{q^{-4}}.
\]

Proceeding by induction, we establish the moments of the stationary distribution:
\[
M(\pi_\alt, 0) = 1, \quad M(\pi_\alt, 2) = 1 + q^{-1}, \quad M(\pi_\alt, 4) = (1+q^2)(1+q^{-1})q^{-1}.
\]
These moment calculations lead to two important results. First, the inner product of $\nu$ with itself:
\[
\langle \nu, \nu \rangle_{\pi_\alt} = \tfrac{q(q-1)}{q+1}.
\]
Second, the total variation norm of $\nu$:
\[
\|\nu\|_{tv} = \tfrac{2\alpha(q)}{(q+1)\eta_1(q)}.
\]

With these preparations, we can now determine the spectral projection of $\delta_0$ onto $V_{q^{-2}}$:
\[
(\delta_0)_{q^{-2}} = q^{-1}(q-1)^{-1} \nu,
\]
which consequently gives:
\[
\|(\delta_0)_{q^{-2}}\|_{tv} = \tfrac{2\alpha(q)}{(q-1)^2(q+1)}.
\]

Turning now to $Q_\alt$, we follow a parallel approach. The corresponding eigenvector is:
\[
\nu' := \pi_\alt' - \tfrac{1}{q+1}(\pi_\alt' \circ q^2) \in V_{q^{-2}}.
\]
Similarly, we identify an element in $V_{q^{-4}}$:
\[
\pi_\alt'-\tfrac{1}{q}(\pi_\alt'\circ q^2)+\tfrac{1}{q(q+1)(q^2+1)}(\pi_\alt'\circ q^4) \in V_{q^{-4}}.
\]

The moment calculations for $Q_\alt$ yield:
\[
M(\pi_\alt', 0) = 1, \quad M(\pi_\alt', 2) = q+1, \quad M(\pi_\alt', 4) = (q+1)(q^2+1).
\]
From these, we derive the key quantities:
\[
\langle \nu', \nu' \rangle_{\pi_\alt'} = \tfrac{q(q-1)}{q+1}, \quad \|\nu'\|_{tv} = \tfrac{\alpha(q) q}{q+1}.
\]

Finally, the spectral projection for $Q_\alt$ satisfies:
\[
(\delta_0)_{q^{-2}} = (q-1)^{-1} \nu', \quad \|(\delta_0)_{q^{-2}}\|_{tv} = \tfrac{\alpha(q) q}{(q-1)(q+1)}.
\]
Then the desired results follows from Theorem \ref{thm:main_thm_spec}.
\end{proof}

\end{document}
